\documentclass[11pt,a4paper]{article}

\usepackage[utf8]{inputenc}
\usepackage[T1]{fontenc}
\usepackage[english]{babel}
\usepackage[margin=2.3cm]{geometry}
\usepackage{amsmath,amssymb}
\usepackage{listings}
\usepackage{verbatim}
\usepackage{xcolor}
\usepackage{hyperref}
\usepackage{booktabs}
\usepackage{float}
\usepackage{enumitem}

\hypersetup{colorlinks=true, linkcolor=blue, urlcolor=blue, citecolor=blue}

\definecolor{lstbg}{rgb}{0.97,0.97,0.97}
\definecolor{lstframe}{rgb}{0.70,0.70,0.70}

\lstset{
    basicstyle=\ttfamily\footnotesize,
    keywordstyle=\color{blue}\bfseries,
    commentstyle=\color{gray}\itshape,
    stringstyle=\color{orange},
    showstringspaces=false,
    breaklines=true,
    frame=single,
    rulecolor=\color{lstframe},
    backgroundcolor=\color{lstbg},
    tabsize=4,
    captionpos=b,
}

\lstdefinestyle{plain}{
    basicstyle=\ttfamily\footnotesize,
    breaklines=true,
    frame=single,
    rulecolor=\color{lstframe},
    backgroundcolor=\color{lstbg},
    keywordstyle={},
    commentstyle={},
    stringstyle={},
}

\title{\textbf{Final Report on the Project ``Solving the Subset Sum
Problem via Quantum Walk Search''}%

\large
A study, implementation analysis and critical comparison of the
MNRS quantum walk algorithm for CSSP.

\url{https://github.com/jacopobellosi/Solving-the-Subset-Sum-Problem-via-Quantum-Walk-Search}
}

\author{Jacopo Bellosi}
\date{April 2026}

\begin{document}
\maketitle

\begin{abstract}
\noindent This report summarises the work carried out on the paper
\emph{``Solving the Subset Sum Problem via Quantum Walk Search''}
(Lancellotti et al., IEEE TC, vol.~75, n.~1, 2026). The paper proposes
an MNRS-style quantum walk on the Johnson graph $J(n,k)$ for the
Constrained Subset Sum Problem (CSSP), with a fully-fledged
\texttt{myQLM} reference implementation. We first analysed the paper
and the source code, then attempted to reproduce its behaviour on a
small toy instance amenable to classical state-vector simulation. The
simulation produced a flat probability distribution
$(\tfrac13,\tfrac13,\tfrac13)$ instead of the expected amplification
of the marked solution. We diagnosed the problem with five
incremental modifications, compared the result against six related
quantum-search papers, and concluded that the failure on toy
instances is structural — the regime in which MNRS provides a
quadratic speed-up does not overlap with the regime in which a
state-vector simulator can run. The report is organised in five
parts: introduction to the paper and the algorithm
(Section~\ref{sec:intro}), description of the reference code and the
observed limits (Section~\ref{sec:code}), critical comparison with
the literature (Section~\ref{sec:compare}), the modifications we
applied (Section~\ref{sec:work}), and final conclusions
(Section~\ref{sec:concl}).
\end{abstract}

\tableofcontents

% =====================================================================
\section{The paper, the problem and the algorithm}
\label{sec:intro}
% =====================================================================

\subsection{The Subset Sum Problem and its constrained variant}

The \emph{Subset Sum Problem} (SSP) is a textbook NP-complete problem:
given a multiset of integers $X = \{x_1,\dots,x_n\}$ and a target
$p\in\mathbb{Z}$, decide whether there exists a subset $S\subseteq X$
such that
\[
\sum_{x_i\in S} x_i \;=\; p .
\]
The \emph{Constrained Subset Sum Problem} (CSSP) studied in
Lancellotti et al.\ adds a cardinality constraint: the searched
subset must contain exactly $k$ elements, $|S|=k$. CSSP is just as
hard as SSP, but the cardinality constraint makes the search space
combinatorial in a controlled way: the candidate subsets are exactly
the $\binom{n}{k}$ vertices of the Johnson graph $J(n,k)$.

The interest in CSSP is threefold:
\begin{itemize}[leftmargin=1.2em]
\item \textbf{Cryptography.} CSSP underlies a number of public-key
schemes (Lyubashevsky-Palacio-Segev 2010, Lagarias-Odlyzko 1983) and
is one of the building blocks used to estimate the cost of the best
quantum attack against post-quantum primitives.
\item \textbf{Combinatorial optimisation.} It models knapsack-like
selection problems in scheduling and resource allocation.
\item \textbf{Complexity theory.} As an NP-complete problem, it is
one of the natural targets to test where quantum algorithms can
deliver a provable speed-up.
\end{itemize}

A brute-force classical search costs $O\!\left(\binom{n}{k}\right)$;
the paper's contribution is a quantum algorithm that achieves the
expected Grover-like \emph{quadratic speed-up}
$O\!\left(\sqrt{\binom{n}{k}}\right)$, but inside the more refined
MNRS quantum walk framework rather than via a direct Grover oracle.

\subsection{The MNRS quantum walk in one page}

The MNRS framework (Magniez-Nayak-Roland-Santha 2007) generalises
Grover's amplitude amplification by replacing the unstructured search
over a database with a \emph{quantum walk} over a graph $G=(V,E)$.
The total cost of MNRS is
\[
\mathrm{Cost}_{\mathrm{MNRS}}
\;=\; T_S \;+\; \frac{1}{\sqrt{\varepsilon}}
\!\left(\frac{T_U}{\sqrt{\delta}} + T_C\right),
\]
where $T_S$ is the setup cost (preparing the stationary state),
$T_U$ is the cost of one walk step, $T_C$ the cost of one oracle
check, $\varepsilon$ the fraction of marked vertices and
$\delta$ the spectral gap of the walk. The quadratic speed-up over a
classical random walk lives in the $1/\sqrt{\delta}$ factor: a small
gap means a slowly-mixing chain, where the quantum walk wins; a
spectral gap of order $1$ collapses the speed-up to a constant.

The four ingredients of the algorithm are sketched below; we will
describe each more precisely in Section~\ref{sec:code}.

\paragraph{(1) Johnson graph $J(n,k)$.} A vertex is a $k$-subset
$S\subseteq\{1,\dots,n\}$; two vertices are connected by an edge iff
they share $k-1$ elements (i.e.\ differ in one element).
$J(n,k)$ has $\binom{n}{k}$ vertices, regular degree $k(n-k)$ and a
spectral gap known in closed form,
\begin{equation}
\label{eq:gap}
\delta \;=\; \frac{n}{k(n-k)},
\end{equation}
valid for $k(n-k)\ge n$ (we will see in Section~\ref{sec:code} that
this regime hypothesis is what makes the toy instance pathological).

\paragraph{(2) Marking oracle $U_f$.} Quantises the indicator
function of the solution: $U_f|S\rangle = (-1)^{f(S)}|S\rangle$, where
$f(S)=1$ iff $\sum_{x\in S}x = p$. In practice this is implemented
by computing the sum of the elements of $S$ in an ancillary register
and comparing it bitwise against the target $p$.

\paragraph{(3) Walk operator $U_W$.} The Szegedy quantisation of one
step of the random walk on $J(n,k)$. It is built as a double
reflection,
\[
U_W \;=\; U_{\mathrm{ref}}(a_{v_j})\,U_{\mathrm{ref}}(a_{v_i})
\;=\; U_U U_U\,U_{\mathrm{ref}}(0)\,U_U U_U\,U_{\mathrm{ref}}(0)\,U_U,
\]
where $U_U$ (the ``update operator'') generates the superposition of
the neighbours of the current vertex.

\paragraph{(4) Approximate reflection via QPE.} The amplitude
amplification step needs a reflection $U_{\mathrm{ref}}(E)$ around
the stationary state $|E\rangle$ of the walk. Since $|E\rangle$ has
eigenphase $0$ under $U_W$ and the other eigenstates have non-zero
eigenphase, a Quantum Phase Estimation (QPE) on $U_W$ acts as an
approximate reflection: it flips the sign of every state whose
estimated phase is nonzero. The standard MNRS precision requirement
fixes the number of QPE qubits to
\begin{equation}
\label{eq:ls}
\ell_s \;=\; \left\lceil \log_2\!\frac{\pi}{2\sqrt{\delta}} \right\rceil .
\end{equation}

The algorithm interleaves $\sqrt{1/(\varepsilon\,\delta)}$
applications of $U_f$ and $U_{\mathrm{ref}}(E)$ to concentrate the
amplitude on the marked subspace. With high probability, a final
measurement of the vertex register returns a marked vertex.

\subsection{What the paper claims and what it does not}

The paper provides:
\begin{itemize}[leftmargin=1.2em]
\item Closed-form analytical estimates of the gate count, depth and
qubit count of every sub-routine, parametrised by $(n,k)$ and the
bit-width $m$ of the input integers.
\item A comparison against a Grover baseline showing that the walk
gives a strictly better T-count for $n\gtrsim 32$ and a strictly
better T-depth for $n\gtrsim 8192$.
\item Two reference open-source implementations, in Qiskit and
\texttt{myQLM}, of all the building blocks (Dicke state preparation,
update operator in two variants, oracle, QPE-based reflection).
\end{itemize}

The paper, however, makes \textbf{no end-to-end claim about the
output distribution} on a concrete instance. The authors validate
the building blocks separately and explicitly state that an
end-to-end execution is not feasible, both on hardware (the algorithm
targets fault-tolerant quantum computers, with error rate
$\sim 10^{-15}$) and on a classical simulator (already the smallest
non-trivial instance $J(4,2)$ requires 38 qubits, well above the
$\sim 30$-qubit ceiling of a state-vector simulator on commodity
hardware). The cryptographic regime targeted by the paper is
$n\ge 25$ for the T-count advantage and $n\ge 2^{13}$ for the
T-depth advantage.

\subsection{What our project set out to do}

Our goal was to:
\begin{enumerate}[leftmargin=1.4em]
\item Read and understand the paper and its reference
\texttt{myQLM} code (\texttt{paper-codes/2025-TC}, mirrored locally
in \texttt{code\_2025TC/}).
\item Run the code end-to-end on the smallest possible toy instance
on \texttt{PyLinalg}, observe the output distribution and check that
the marked solution is amplified.
\item Diagnose any discrepancy and document each modification needed
to bring the simulation in line with the theory.
\item Compare the result with related works in the literature
(folder \texttt{documentation/paper\_comparison/}) and draw
conclusions on the practical reproducibility of the algorithm.
\end{enumerate}

The chosen toy instance is $n=3$, $k=1$,
$\mathtt{values}=[1,2,3]$, $\mathtt{target\_sum}=3$. With
$N=\binom{3}{1}=3$ admissible vertices and the unique marked
solution $S=\{3\}$, an amplifying algorithm should produce a
distribution close to $(0,0,1)$ on the vertex register. Instead,
the reference code produces $(\tfrac13,\tfrac13,\tfrac13)$ — the
initial Dicke distribution, untouched.

% =====================================================================
\section{The reference code and the observed limits}
\label{sec:code}
% =====================================================================

\subsection{Anatomy of the \texttt{myQLM} implementation}

The reference code is structured around the file
\texttt{code\_2025TC/cssp.py}, which orchestrates a number of
sub-routines from the \texttt{qatext} package (Dicke state, VBE
adder, Cuccaro arithmetic, Bartschi Dicke generator, sliding sort
arrays, etc.). The circuit operates on the following named
quantum registers (Table~\ref{tab:registers}).

\begin{table}[H]
\centering
\footnotesize
\begin{tabular}{@{}l l l@{}}
\toprule
\textbf{Register} & \textbf{Width} & \textbf{Role} \\
\midrule
$\sigma$ (\textsf{dicke})    & $n$ qubits     & Dicke state $|D^n_k\rangle$ on the indicator bits of $S$ \\
$S_1, S_0$                   & $k\cdot m$ and $(n-k)\cdot m$ & Sorted values inside / outside the current vertex \\
$T_1, T_0$                   & idem            & Sorted values for the neighbour vertex (used by $U_W$)\\
$\alpha_1, \alpha_0$         & $m$ qubits      & Element being inserted/removed during update \\
$\omega_1, \omega_0$         & $k$ and $n-k$ qubits & Bartschi $W$-states selecting which element to move \\
\textsf{qpe\_s}              & $\ell_s$ qubits & Phase register for the approximate reflection \\
\textsf{sum\_reg}            & $\lceil\log_2 k\rceil + m$ & Accumulator for the oracle's sum \\
\bottomrule
\end{tabular}
\caption{The eight register groups used by \texttt{cssp.py}. For
the toy instance $(n,k,m)=(3,1,2)$, the total qubit count is 28
(rising to 30 after Modification~4 forces $\ell_s\ge 3$).}
\label{tab:registers}
\end{table}

The circuit is built incrementally with \texttt{myQLM}'s
\texttt{Program}/\texttt{compute()}/\texttt{uncompute()} mechanism.
The high-level structure of one MNRS iteration is
\[
|\psi_0\rangle
\xrightarrow{\text{Dicke + VBE}} |\psi_1\rangle
\xrightarrow{U_U} |\psi_2\rangle
\xrightarrow{U_f} |\psi_3\rangle
\xrightarrow{U_{\mathrm{ref}}(E)} |\psi_4\rangle ,
\]
with $U_{\mathrm{ref}}(E)$ realised as a QPE-based block sandwiched
between Reflection~A and Reflection~B (the two Householder
reflections that compose into $U_W$). The paper offers two
implementations of $U_U$:
\begin{itemize}[leftmargin=1.2em]
\item \emph{Low-width}: sequential, fewer ancillas, depth $\sim 10nm$.
The \texttt{insert}/\texttt{delete} pair is reversible by
construction.
\item \emph{Low-depth}: parallel, more ancillas, depth
$O(\log m)$. The \texttt{insert\_ld} routine in
\texttt{sliding\_sort\_array.py} is flagged in-source as
``not properly reversible (garbage in ancillas)''.
\end{itemize}
This distinction will be relevant in Section~\ref{sec:work}
(Modification~3).

\subsection{Running the code as-is on the toy instance}

A direct invocation
\texttt{python code\_2025TC/cssp.py True} on the toy instance
$(n,k)=(3,1)$, \texttt{values}=[1,2,3], \texttt{target}=3 builds a
circuit with 28 qubits and about 2547 gates and submits it to
\texttt{PyLinalg}. The simulation completes without errors after
roughly 11 hours of single-core CPU time and produces the output
shown in the left column of Table~\ref{tab:baseline-final}.

\begin{table}[H]
\centering
\footnotesize
\begin{tabular}{@{}p{0.48\linewidth} p{0.48\linewidth}@{}}
\toprule
\textbf{(a) Direct run of \texttt{cssp.py}} &
\textbf{(b) Run with intermediate checkpoints} \\
\midrule
\begin{minipage}[t]{\linewidth}
\begin{lstlisting}[style=plain, basicstyle=\ttfamily\scriptsize]
n=3, k=1, values=[1,2,3], target=3
{'nbqbits': 28, 'gate_size': 2547}

Final measurement on S_1:
  0.33333  |1>
  0.33333  |2>
  0.33333  |3>

(runtime: ~11h CPU on a single core)
\end{lstlisting}
\end{minipage}
&
\begin{minipage}[t]{\linewidth}
\begin{lstlisting}[style=plain, basicstyle=\ttfamily\scriptsize]
Ckpt 1  Dicke + VBE:
  6 states, p=0.167, ampl=+0.408
  uniform over the 3 subsets

Ckpt 2  Edge superposition (U_U):
  12 states, p=0.083, ampl=+0.289
  uniform over all edges of J(3,1)

Ckpt 3  Post-Oracle:
  ampl=+0.289 on Sum=2
  ampl=-0.289 on Sum=3  <-- phase flip OK

Ckpt 4-5  Post-QPE / Post-Diffusion:
  signs scrambled, no concentration

Ckpt 6  End of iteration:
  flat prob on S_1
\end{lstlisting}
\end{minipage}\\
\bottomrule
\end{tabular}
\caption{Output of the original code. The final distribution on
$S_1$ is identical to the initial Dicke distribution. The oracle
correctly applies the phase flip on the marked states; the
QPE-based reflection scatters the signs instead of rotating them
into amplitude.}
\label{tab:baseline-final}
\end{table}

The result is unexpected: the algorithm should amplify the marked
solution $S=\{3\}$, producing a distribution close to $(0,0,1)$.
Instead, the final distribution on $S_1$ is exactly the initial
Dicke distribution $|D^3_1\rangle$, untouched by the iterations.

\subsection{The diagnostic checkpoints}

The same circuit was instrumented with checkpoints at six key points
of one MNRS iteration (the file
\texttt{code\_2025TC/cssp\_with\_checkpoint.py} produces the trace).
The right column of Table~\ref{tab:baseline-final} summarises what
each checkpoint shows. The crucial observation is at Checkpoint~3
(post-oracle): all 12 amplitudes are
$\pm 0.2887 = \pm 1/\sqrt{12}$, with the correct phase flip on the
states whose sum is $3$ (the marked sum). \emph{The oracle is
working}.

The breakdown happens between Checkpoint~3 and Checkpoint~5: the
QPE-based reflection scrambles the signs instead of rotating them
into amplitude. By Checkpoint~6 the marginal distribution on $S_1$
is uniform.

This is the empirical fingerprint that motivated the work described
in Section~\ref{sec:work}.

\subsection{Three flavours of ``why is the output flat''}

Before discussing the modifications, it is useful to separate three
\emph{independent} causes that all manifest themselves as a flat
output distribution. Conflating them leads to chasing the wrong fix.

\paragraph{(C1) Physical limit of state-vector simulation.}
Already on $J(3,1)$ the circuit allocates 28 qubits, which means a
state vector of $2^{28}\approx 2.7\cdot 10^8$ complex amplitudes,
touched by $\sim 2500$ gates per iteration. A genuine MNRS regime
($\delta < 1$) would require $J(5,2)$ or $J(6,3)$, pushing the
qubit count above 35 and the state-vector size by a factor of
$\sim 128$. The simulator simply cannot follow.

\paragraph{(C2) Structural limit of MNRS on trivial instances.}
$J(3,1)$ is the triangle $K_3$. A walker on $K_3$ mixes in one
classical step; the quantum walk operator has only three eigenphases
$\{0,\pm \arccos(-\tfrac12)\}$ and the formal value
$\delta=n/(k(n-k))=1.5$ is outside the physically meaningful
range $(0,1]$. With $1/\sqrt{\delta}=O(1)$ the MNRS quadratic
speed-up collapses to a constant: there is no slow direction for
QPE to discriminate $|E\rangle$ from the rest of the spectrum.

\paragraph{(C3) Genuine bugs in the reference code.}
The chronology of the modifications (Section~\ref{sec:work})
identified three bona-fide bugs in \texttt{cssp.py}: an asymmetric
Reflection~B, a destructive cleanup block of the QPE ancillas and a
non-reversible \texttt{insert\_ld} composed back-to-back with
\texttt{delete}. These are real bugs, but their effect on the output
is invisible on $J(3,1)$ because (C2) leaves no amplification
window where the bugs would matter.

The three causes are not in tension: they are stacked. Even if all
the bugs were fixed (they are now), (C2) would still leave the
output flat on $J(3,1)$, and (C1) would prevent us from running the
fixed code on an instance where (C2) does not bite.

% =====================================================================
\section{Comparison with the literature}
\label{sec:compare}
% =====================================================================

To put the reproducibility difficulties of Lancellotti et al.\ in
perspective, we surveyed six related papers stored in
\texttt{documentation/paper\_comparison/}. The selection covers the
two natural baselines (a direct quantum algorithm for SSP on toy
instances; a Grover-based quantum attack on a related cryptographic
problem) and four works on quantum walks where classical simulation
or hardware execution \emph{did} succeed. The detailed per-paper
notes are summarised in Table~\ref{tab:papers}; the rest of this
section explains why those works are not refutations of our
findings, but rather confirmations.

\begin{table}[H]
\centering
\footnotesize
\setlength{\tabcolsep}{4pt}
\begin{tabular}{@{}p{0.18\linewidth} p{0.18\linewidth} p{0.13\linewidth} p{0.18\linewidth} p{0.25\linewidth}@{}}
\toprule
\textbf{Reference} & \textbf{Algorithm} & \textbf{Instance simulated} & \textbf{Validation platform} & \textbf{Why it does not address our regime}\\
\midrule
Zheng et al., 2022 (\texttt{s11432-021-3334-1}) &
QPE + Amplitude Amplification for SSP &
$|S|{=}4$, $w{=}12$, $10$ qubits &
Qiskit + IBM Q hardware (Santiago, Bogota) &
Sidesteps quantum walk entirely; oracle marks the sum directly with no Johnson graph \\
\midrule
Sato et al., 2024 (\texttt{2405.07129v3}) &
Two-step Grover for TSP, HOBO encoding &
$n{=}3,4$ cities, $6$/$12$ qubits &
Qiskit Aer (state-vector) &
Compact encoding bypasses spectral-gap issues; no walk \\
\midrule
Perriello et al., 2021 (\texttt{ISD.pdf}) &
Grover for ISD on Prange formulation &
crypto sizes ($\sim 2^{32}$ qubits) &
Atos QLM, sub-component validation only &
Explicitly admits no end-to-end simulation; only resource estimates \\
\midrule
Razzoli et al., 2024 (\texttt{entropy-26-00313}) &
DTQW on $2^n$-cycles, no marking &
$N=4,8$, hardware run &
IBM \textsf{ibm\_cairo} (27 qubits) &
No oracle, no reflection; bypasses the spectral-gap problem entirely \\
\midrule
da Silva et al., 2024 (\texttt{entropy-26-00668}) &
Grover on database of 16 colours &
$5$ qubits &
PyLinalg + NoisyQProc &
Standard Grover on $N=16$, well inside the regime $1/\sqrt N$ \\
\midrule
Portugal \& Moqadam, 2025 (\texttt{entropy-27-00454-v2}) &
CTQW search on $K_N$, $K_{n,n}$, $Q_n$ &
$K_8$, $K_{32,32}$, $Q_{10}$ &
Qiskit + Hiperwalk &
Uses analytical eigenvectors of highly symmetric graphs; valid only asymptotically on $Q_n$ \\
\midrule
Wing-Bocanegra et al., 2025 (\texttt{s41598-025-87304-0}) &
Single-step DTQW on $K_{2^n}$ for QAOA init &
$K_4$ to $K_{64}$ &
Qiskit + IBM \textsf{ibmq\_manila} &
A single walk step removes the iteration-tuning problem and the spectral-gap dependence \\
\bottomrule
\end{tabular}
\caption{Comparison of related quantum-search papers. None of them
implements the MNRS framework on a Johnson graph at the scale
targeted by Lancellotti et al.; all of them succeed by avoiding the
specific obstacle that defeats the toy-instance simulation of CSSP.}
\label{tab:papers}
\end{table}

\subsection{The crucial common pattern: avoiding the regime
boundary}

Reading these works in parallel, one pattern emerges: every paper
that successfully validates a quantum algorithm in simulation does
so by sidestepping the precise constraint that breaks our toy
simulation. We can group the avoidance strategies in three
families.

\paragraph{Family A — replace MNRS by a different algorithm.}
Zheng et al.~(2022) is the closest analogue to Lancellotti et al.:
they too solve SSP, on a toy instance ($n{=}4$, $w{=}12$), with
\emph{Quantum Phase Estimation followed by Amplitude Amplification}
on a direct sum oracle, no Johnson graph and no walk. The simulation
succeeds at 94\% and even runs on real NISQ hardware at 75\%.
The price is methodological: there is no Markov-chain interpretation
and the asymptotic complexity is the standard Grover bound, not the
walk speed-up. This is exactly the ``viable alternative \#1'' we
identified at the end of Section~\ref{sec:work}: drop MNRS entirely.

\paragraph{Family B — keep walks but pick easy graphs.}
Razzoli et al.~(2024), Portugal \& Moqadam~(2025) and
Wing-Bocanegra et al.~(2025) all run quantum walks in simulation,
but on grafts whose spectrum is either trivial ($K_N$, $K_{n,n}$
have a two-eigenvalue spectrum), available in closed form
(Hadamard coin on cycles), or sidestepped by running only one step
(Wing-Bocanegra). None of them uses a Johnson graph $J(n,k)$ with
$k\ge 2$, where the spectrum has $k+1$ distinct eigenvalues whose
ratios depend on $n,k$ in a non-trivial way. Equally important,
none of them uses the MNRS \emph{approximate reflection} via QPE:
they either use no reflection at all (free walks), or implement a
direct Grover-style reflection.

\paragraph{Family C — give up on end-to-end simulation.}
Perriello et al.~(2021) is the most candid: working on Information
Set Decoding at cryptographic sizes, they explicitly state that
end-to-end simulation is impossible and only validate the
sub-components (oracle, diffusion). The published numbers are
\emph{resource estimates}, not measured probabilities. This is
methodologically the same trajectory taken by Lancellotti et al.\
themselves: validate the building blocks separately, present
analytical complexity figures, and renounce the end-to-end output
distribution.

\subsection{What the comparison tells us}

The takeaway is unambiguous. The pieces of the MNRS pipeline that
make it powerful at cryptographic scale — the Johnson graph
encoding, the approximate reflection via QPE, the spectral-gap
dependence — are precisely the pieces that make it impossible to
validate at toy scale. The literature's working strategies all
correspond to relaxing one of those pieces:
\begin{itemize}[leftmargin=1.2em]
\item Zheng et al.: drop the walk; keep QPE + Grover.
\item Razzoli et al., Portugal \& Moqadam, Wing-Bocanegra: keep
the walk but on graphs with trivial or analytic spectrum, and
without an MNRS reflection.
\item Perriello et al.: keep the algorithm but never simulate it
end-to-end.
\end{itemize}

The combination ``MNRS + Johnson graph + small instance + classical
simulation + amplitude check'' is, to the best of our literature
survey, an empty intersection. Our finding that the toy run on
$J(3,1)$ produces a flat distribution is consistent with this
pattern: not a bug in the paper, but a structural mismatch between
the regime where MNRS is meaningful and the regime where a
state-vector simulator can run.

% =====================================================================
\section{What we did}
\label{sec:work}
% =====================================================================

This section summarises the five chronological modifications applied
to \texttt{cssp.py}. Each modification was paired with a fresh
simulation log, stored in \texttt{documentation/quantum\_simulation/}
(numbered 01--13). The ``before/after'' diffs and gate counts are
reported in detail in
\texttt{documentation/modification\_report/modification\_report.tex};
here we only report the rationale, the effect and the verdict on
each modification's necessity.

\subsection{Modification 1 — Symmetric Reflection A and B on the
full coin (17 March 2026)}

The MNRS walk operator is a double reflection
$U_W = U_U U_U R_0 U_U U_U R_0 U_U$ where $R_0=2|0_n\rangle\langle 0_n|-I$
acts on the \emph{whole} coin register $\omega = (\omega_1,\omega_0)$
of $n$ qubits. In the original code, both Reflection~A and
Reflection~B implemented $R_0$ asymmetrically: Reflection~A
controlled the $Z$ flip on $\omega_1$ only (\texttt{wstate\_ones},
$k$ qubits), and Reflection~B did the analogous thing on $\omega_0$,
with a typo in the second X-sandwich (\texttt{range(k)} instead of
\texttt{range(n-k)}). The fix promotes both reflections to
\texttt{Z.ctrl(n)} on the full coin and corrects the loop bound.

\begin{lstlisting}[language=Python]
# After fix (extract): symmetric Reflection B on the full coin
prw.apply(qrw_update.dag(),
          node_t_ones, node_t_zeros,
          node_s_ones, node_s_zeros,
          alpha_ones, alpha_zeros,
          wstate_ones, wstate_zeros)
for j in range(k):     prw.apply(X, wstate_ones[j])
for j in range(n - k): prw.apply(X, wstate_zeros[j])
prw.apply(Z.ctrl(n), qpe_s[qw_iter],
          wstate_ones, wstate_zeros)              # control on full coin
for j in range(k):     prw.apply(X, wstate_ones[j])
for j in range(n - k): prw.apply(X, wstate_zeros[j])
prw.apply(qrw_update,
          node_t_ones, node_t_zeros,
          node_s_ones, node_s_zeros,
          alpha_ones, alpha_zeros,
          wstate_ones, wstate_zeros)
\end{lstlisting}

\noindent\textbf{Effect.} Gate count rises from 2547 to 2581, qubit
count unchanged (28). Distribution still flat on $J(3,1)$.\\
\noindent\textbf{Verdict.} Genuine bug. Necessary for the algorithm
to be the MNRS walk on the paper. Invisible on $J(3,1)$ because of
cause (C2).

\subsection{Modification 2 — Removal of the $\alpha/\omega$ cleanup
block (20--23 March 2026)}

Inside the QPE block, after the $\ell_s$ controlled walk steps, the
original code performed a manual ``cleanup'' of the ancillary
registers $\alpha$ and $\omega$:
\texttt{copy\_register(m).ctrl()} gates followed by
\texttt{generate(...)}. The intent was to uncompute the ancillas to
preserve reversibility. The mistake is that at that point $\alpha$
and $\omega$ are not in $|0\rangle$: they are entangled with the
visited vertices. Calling \texttt{generate} on a non-zero state
applies a non-trivial unitary that corrupts the entanglement
instead of restoring it. The correct uncompute is already provided
by \texttt{prw.uncompute()} at the end of the \texttt{compute()}
block.

\noindent\textbf{Effect.} Gate count drops from 2547 to 1359
(\textbf{$-47\%$}), qubit count unchanged. Distribution still flat
on $J(3,1)$.\\
\noindent\textbf{Verdict.} Genuine bug. The cleanup was both
redundant and destructive. On instances with a non-trivial
amplification window this would have introduced spurious phases
across the QFT, making the algorithm \emph{anti-amplify}. Invisible
on $J(3,1)$ for the same reason as Modification~1.

\subsection{Modification 3 — Use \texttt{insert\_lw} instead of
\texttt{insert\_ld} (22 March 2026)}

The \texttt{delete} routine in
\texttt{qatext/qroutines/datastructure/sliding\_sort\_array.py} was
internally calling \texttt{insert\_ld(...).dag()}. A reversibility
sanity check failed: \texttt{delete} composed back-to-back with
\texttt{insert\_ld} did not act as the identity on a small
non-ancilla input (a comment in source already flagged
\texttt{insert\_ld} as ``not properly reversible — garbage in
ancillas''). Since the walk operator $U_W$ composes many
\texttt{insert}/\texttt{delete} calls, garbage propagates and
degrades QPE. The fix routes \texttt{delete} through
\texttt{insert\_lw(...).dag()}, which passes the sanity check.

\noindent\textbf{Effect.} Gate count 1359 $\to$ 1366 (negligible
overhead, the low-width variant has slightly more sequential CNOTs).
Distribution still flat on $J(3,1)$.\\
\noindent\textbf{Verdict.} Structural bug. The walk operator was
not unitary on the ancillary subspace before this fix. Necessary
for the algorithm to be \emph{well-defined}, even though invisible
on $J(3,1)$.

\subsection{Modification 4 — Clamp $\delta\le 1$ and force
$\ell_s\ge 3$ (27 March 2026)}

Equation~\eqref{eq:gap}, $\delta = n/(k(n-k))$, is derived under the
hypothesis $k(n-k)\ge n$. On $(n,k)=(3,1)$ it returns $\delta=1.5$,
a numerical artifact (a true spectral gap of a stochastic matrix
lies in $(0,1]$). Plugging this back into Eq.~\eqref{eq:ls} gives
$\ell_s = 1$: a single QPE qubit can only distinguish phase $0$ from
phase $\pi$, which is not enough to tell $|E\rangle$ apart from any
other eigenstate of $U_W$. As a defensive patch we clamp $\delta$ to
$1$ and force $\ell_s$ to at least $3$.

\begin{lstlisting}[language=Python]
delta  = min(n / (k * (n - k)), 1.0)            # clamp to physical range
len_s  = max(int(np.ceil(np.log2(np.pi / (2 * np.sqrt(delta))))), 3)
\end{lstlisting}

\noindent\textbf{Effect.} Qubit count rises from 28 to 30; gate count
unchanged (1366). Distribution still flat on $J(3,1)$.\\
\noindent\textbf{Verdict.} Defensive patch, not an algorithmic fix.
It cannot create a spectral gap that the graph does not have. On
graphs with $\delta<1$ it never triggers, so it is fully backward
compatible. On $J(3,1)$ it makes the QPE block at least
\emph{well-defined}, which is a precondition for any meaningful
behaviour.

\subsection{Modification 5 — Checkpoint debugger for amplitude
inspection (17 April 2026)}

The instrumented copy \texttt{cssp\_with\_checkpoint.py} simulates
the full state at six points of the circuit and prints, for each
state with $p>5\cdot 10^{-4}$, the probability and the complex
amplitude, decoded into the named registers. It is not an
algorithmic change — it is a diagnostic tool — but it was decisive
for the project: it pinpointed where the amplification is lost.

\begin{lstlisting}[style=plain]
Ckpt 3 (Post-Oracle):
  Prob 0.0833 | Ampl +0.2887 | S_1:0 | Sum:2     (non-marked)
  Prob 0.0833 | Ampl -0.2887 | S_1:1 | Sum:3     (marked: phase flip OK)

Ckpt 5 (Post-Diffusion):
  Prob 0.0417 | Ampl +/-0.2041 | scattered signs
  (no concentration on Sum:3)
\end{lstlisting}

\noindent\textbf{Effect.} Empirical confirmation that the
oracle is correct (phase flip at Checkpoint~3) and that the failure
sits inside the QPE-based reflection (between Checkpoints 3 and 5).
This excludes any hypothetical further bug in the oracle, in the
Dicke preparation or in the VBE adder.\\
\noindent\textbf{Verdict.} Indispensable observability tool. Without
it, the diagnosis ``the regime is wrong, not the code'' would have
remained a conjecture.

\subsection{Summary of the modifications}

\begin{table}[H]
\centering
\footnotesize
\begin{tabular}{@{}l l l l l@{}}
\toprule
\textbf{Mod} & \textbf{Date} & \textbf{Type} & \textbf{Effect on gates} & \textbf{Effect on distribution}\\
\midrule
1 & 17 Mar & Genuine bug fix    & 2547 $\to$ 2581 & flat \\
2 & 20 Mar & Genuine bug fix    & 2547 $\to$ 1359 ($-47\%$) & flat \\
3 & 22 Mar & Structural fix     & 1359 $\to$ 1366 & flat \\
4 & 27 Mar & Defensive patch    & 1366 (28$\to$30 qubits) & flat \\
5 & 17 Apr & Diagnostic tool    & no change & flat (now \emph{understood}) \\
\bottomrule
\end{tabular}
\caption{The five modifications, in chronological order. Mods 1--3
are real bug fixes that would matter on a genuine MNRS instance;
Mod 4 is a numerical safeguard; Mod 5 is the observability layer
that closed the diagnosis. The output distribution remains flat on
$J(3,1)$ because of cause (C2), independent of the bugs.}
\label{tab:mods}
\end{table}

\subsection{Why no further code-level fix can rescue $J(3,1)$}

Combining the spectral, mixing and precision arguments:
\begin{itemize}[leftmargin=1.2em]
\item \emph{Spectral.} $J(3,1)\equiv K_3$ has only three eigenphases
under $U_W$. The formal $\delta=1.5$ is outside $(0,1]$ and the
honest gap (when computed correctly from $|\lambda_2|$ of the
walk on $K_3$) is at most $\tfrac12$, far from a separation
that QPE with 1--3 qubits can resolve.
\item \emph{Mixing.} A walker on $K_3$ mixes in one classical step.
$1/\sqrt{\delta} = O(1)$ collapses the MNRS speed-up to a constant,
and the cost formula
$1/\sqrt{\varepsilon\,\delta}$ predicts roughly $2$ amplification
iterations: nothing for the algorithm to converge to.
\item \emph{Precision.} The MNRS condition $2^{\ell_s}\ge\pi/(2\sqrt\delta)$
demands $\ell_s$ large enough to resolve the gap. With
$\delta\ge 1$ formally, the formula returns $\ell_s=1$ and even
forcing $\ell_s=3$ provides $8$ phase bins on $3$ eigenphases — a
collision-prone QPE that acts as a sign scrambler, not as a
reflection.
\end{itemize}

In Table~\ref{tab:regimes} we map the toy instances against the
parameters that govern the MNRS regime, and locate the
``simulability window''.

\begin{table}[H]
\centering
\footnotesize
\begin{tabular}{@{}l l l l l l@{}}
\toprule
\textbf{Instance} & $N=\binom{n}{k}$ & $\delta$ formula & $\ell_s$ formula & Approx.\ qubits & Verdict\\
\midrule
$J(3,1)\equiv K_3$ & 3   & 1.5 (clamped 1) & 1 (forced 3) & 28--30 & \emph{Outside MNRS regime}\\
$J(4,2)$            & 6   & 1.0 (borderline) & 2 (forced 3) & $\sim 33$ & Borderline, no margin\\
$J(5,2)$            & 10  & 0.83             & 2            & $\sim 37$ & Genuine MNRS, but $>$30 qubits\\
$J(6,3)$            & 20  & 0.67             & 2            & $\sim 45$ & Genuine MNRS, simulator-blocked\\
$J(7,3)$            & 35  & 0.58             & 3            & $\sim 52$ & Genuine MNRS, simulator-blocked\\
\bottomrule
\end{tabular}
\caption{The simulability window for the paper-faithful MNRS
implementation. Instances small enough to simulate are outside the
MNRS regime; instances inside the regime exceed the qubit budget of
a state-vector simulator. The intersection is empty.}
\label{tab:regimes}
\end{table}

% =====================================================================
\section{Conclusions}
\label{sec:concl}
% =====================================================================

\subsection{What we learned about the paper}

The paper of Lancellotti et al.\ proposes a careful and detailed
MNRS implementation for CSSP. Reading the paper end-to-end and
implementing the diagnostic checkpoints clarified two facts that
are not explicit in the published text:
\begin{enumerate}[leftmargin=1.4em]
\item The reference \texttt{myQLM} code contained three bugs of
varying severity (Modifications 1--3). The first two are
load-bearing on any instance with a real amplification window; the
third breaks unitarity at the level of the update operator. None
were detectable in the paper because the authors do not run the
algorithm end-to-end.
\item The authors' resource estimates and complexity figures
already place the paper well outside the regime of any classical
simulator: the smallest non-trivial instance, $J(4,2)$, is reported
in their own Table~III as needing 38 qubits, above the practical
ceiling. Our toy run on $J(3,1)$ was only made possible by choosing
an instance \emph{below} the regime in which the algorithm is
designed to work.
\end{enumerate}

\subsection{What we learned about the simulation}

The flat output distribution on $J(3,1)$ is the result of three
independent causes (C1--C3) that overlap on the same observational
fingerprint. The diagnostic checkpoints (Modification~5) make it
possible to separate them empirically: the oracle works (phase flip
at Checkpoint~3), the Dicke preparation works (uniform amplitudes
at Checkpoint~1), the failure sits in the QPE-based reflection
between Checkpoints 3 and 5. After the four code-level fixes that
remove the genuine bugs (C3), the residual flat distribution is
entirely the structural mismatch (C2): $J(3,1)\equiv K_3$ does not
have a spectral gap that MNRS can exploit, and no amount of tuning
inside the reflection block can create one.

\subsection{What we learned from the literature comparison}

Across the seven related papers we surveyed
(Section~\ref{sec:compare}, Table~\ref{tab:papers}), no work
implements the combination ``MNRS framework + Johnson graph + toy
instance + classical simulation + amplitude check''. The successful
simulations either (a) replace MNRS by a different quantum
algorithm (Zheng et al.\ use QPE + Amplitude Amplification on a
direct sum oracle), (b) keep walks but pick graphs with trivial or
analytic spectrum (cycles, complete graphs, bipartite graphs,
hypercubes), or (c) restrict themselves to resource estimates and
sub-component validation, exactly as Lancellotti et al.\ do. Our
inability to validate the algorithm on $J(3,1)$ is therefore
\emph{consistent with the literature}, not anomalous: the literature
itself documents no successful MNRS-on-Johnson simulation, and the
working alternatives are precisely the ones we identify as viable
in the modification report.

\subsection{Viable paths forward}

If the goal is to demonstrate amplification on a laptop simulator,
two pragmatic alternatives remain.

\paragraph{Path A — Direct Grover on a compact vertex register.}
For $N=\binom{n}{k}$ in the low tens, index the admissible subsets
on $\lceil\log_2 N\rceil$ qubits, use an oracle that checks
$\sum_{i\in S}x_i = p$, and apply
$\lfloor\frac{\pi}{4}\sqrt{N}\rfloor$ Grover iterations with the
standard reflection around the uniform superposition. This works
where MNRS fails on toy instances, because the diffusion is exact
and does not rely on a QPE block that needs $\ell_s$ qubits of
precision. The cost is methodological: this is no longer the
walk-based algorithm of Lancellotti et al., and the asymptotic
complexity is the standard Grover bound, not the walk speed-up.
This is the path taken by Zheng et al.~(2022).

\paragraph{Path B — Stripped-down MNRS with a compact encoding.}
For a fixed instance like $(5,2)$ or $(6,3)$, drop the Dicke
register, drop the double cover $S_0$ and the $T$ copies, and
initialise $|S\rangle$ directly on a single register of
$\lceil\log_2 N\rceil$ qubits. The QPE block can be replaced by an
explicit reflection built from the analytical eigenvectors of the
walk operator on small Johnson graphs. The resulting circuit uses
roughly 8--14 qubits and is simulable on \texttt{PyLinalg} in
minutes. This is no longer a faithful reproduction of the paper but
a stripped-down variant whose purpose is to demonstrate the
amplification \emph{qualitatively} while keeping the walk structure
intact.

\subsection{Closing remark}

The MNRS pipeline of Lancellotti et al.\ is algorithmically correct
once the two localised bugs in Modifications~1 and~2 (and the
structural fix in Modification~3) are removed; the patched
implementation lives in \texttt{code\_updated/}. It cannot, however,
be executed in a meaningful way on the toy instances $J(3,1)$ and
$J(4,2)$: on those instances $1/\sqrt{\delta}=O(1)$, so MNRS still
runs but produces no amplification window, and $\ell_s$ is too small
for the approximate reflection to be resolved. The paper itself is
explicit about this: its complexity analysis targets the
cryptographic scale ($n\ge 25$ for lower T-count, $n\ge 2^{13}$ for
lower T-depth) on fault-tolerant hardware, and it notes that the
$n\log_2 n$ qubit overhead ``significantly limits the feasibility of
classical simulation''. The two realistic ways forward are to
replace MNRS with a direct Grover search on a compact vertex
register, or to build a stripped-down MNRS variant on a single
register and an exact reflection — keeping the walk structure but
giving up the one-to-one correspondence with the paper.

\bigskip

\noindent\textbf{Repository.}
\url{https://github.com/jacopobellosi/Solving-the-Subset-Sum-Problem-via-Quantum-Walk-Search}

\noindent\textbf{Companion documents.}
\begin{itemize}[leftmargin=1.2em,nosep]
\item \texttt{documentation/modification\_report/modification\_report.tex}
— full diff-level chronology of the five modifications.
\item \texttt{documentation/quantum\_simulation/01--13.txt}
— per-modification simulation logs.
\item \texttt{documentation/paper\_comparison/}
— PDFs of the seven surveyed related works.
\item \texttt{code\_updated/cssp.py}
— patched implementation with all five modifications applied.
\end{itemize}

\end{document}
